% KKchemstruct 取扱説明書（日本語版）
%   lualatex KKchemstruct-manual-ja.tex   （2回）
% \KKchemmark と \KKhbond の位置決めのため、2回コンパイルしてください。
\documentclass[paper=a4paper,fontsize=10pt]{jlreq}
\usepackage{KKchemstruct}

\title{KKchemstruct 取扱説明書}
\author{KKTeX}
\date{2026年8月2日\quad v1.0.0}

\begin{document}
\maketitle
\tableofcontents

\section{概要}

\textsf{KKchemstruct} は、高校化学（有機・無機）の構造式を組むためのマクロ集です。
描画そのものは \textsf{chemfig} が行い、本パッケージは
参考書の図版の見えに合うよう寸法・線幅・向きを設定したうえで、
置換ベンゼン・示性式・反応式・注釈を短い記法で書けるようにします。

寸法はすべて \verb|em| 単位で定義してあるので、本文の文字サイズを変えれば
図もそれに追随します。判型ごとに数値を書き直す必要はありません。

\subsection{必要なもの}

\begin{itemize}
\item TeX エンジン: LuaLaTeX、または pLaTeX + dvipdfmx
\item 依存パッケージ: \textsf{chemfig}、\textsf{etoolbox}、\textsf{xkeyval}
      （いずれも TeX Live に収録）
\end{itemize}

\subsection{読み込み}

\verb|KKchemstruct.sty| を TeX のサーチパス上か文書と同じディレクトリに置き、

\begin{verbatim}
\usepackage{KKchemstruct}
\end{verbatim}

と書きます。オプションはありません。

\subsection{2回のコンパイルが要る場合}

\verb|\KKchemmark| と \verb|\KKhbond| は TikZ のノード参照を使うため、
位置が定まるまでに2回のコンパイルが必要です。1回目は枠や点線が
ずれた場所に出ますが、2回目で正しい位置に落ち着きます。

\section{ベンゼン環}

\subsection{\texttt{\char`\\KKbenzene}}

置換ベンゼンは \verb|\KKbenzene| で書きます。置換位置は真上を 1 として
時計回りに 1〜6 です。

\begin{verbatim}
\KKbenzene[2=OH,3=COOH]
\end{verbatim}

\begin{center}
\KKbenzene[2=OH,3=COOH]\qquad
\KKbenzene[1=COOH,4=COOH]\qquad
\KKbenzene[]
\end{center}

置換基は数式モードで組まれるので、\verb|COOCH_3| のように添字を
そのまま書けます。空の位置は書かなければ枝ごと出ません。

\subsection{オプション}

\begin{description}
\item[\texttt{1}〜\texttt{6}] 各位置の置換基。
\item[\texttt{angle}] 環全体の回転角（度）。既定は回転なし（頂点が上下）。
\item[\texttt{sep}] この環だけの一辺の長さ（原子中心間距離）。
\item[\texttt{kekule=b}] 二重結合の内側短線を裏返し、右上・下・左上の3辺に入れます。
      置換基と短線がぶつかって見苦しいときに使います。
\end{description}

\begin{verbatim}
\KKbenzene[1=OH,kekule=b]      \KKbenzene[2=OH,angle=30,sep=1.0em]
\end{verbatim}

\begin{center}
\KKbenzene[1=OH,kekule=b]\qquad
\KKbenzene[2=OH,angle=30,sep=1.0em]
\end{center}

\subsection{本文中に置く横向きの環}

段落の途中に環を流し込むときは、上下の辺が水平な（flat-top の）形を使います。
左右の頂点で結合するので、鎖の途中に自然に収まります。

\begin{description}
\item[\texttt{\char`\\KKphenylring}] 無置換の環だけ。
\item[\texttt{\char`\\KKphenyl\{...\}}] 環の右へ置換基を出します。
\item[\texttt{\char`\\KKphenylene\{...\}}] 左にも結合の切り株を出し、鎖の途中に挟めます。
\end{description}

本文の中に \KKphenylring{} を置いても行送りは崩れません。
\KKphenyl{COOH} のように右へ置換基が出せますし、
鎖の途中に \ck{CH_3(CH_2)_{11}}\KKphenylene{SO_3H} と挟むこともできます。

ベースラインからの持ち上げ量は \verb|\KKchemsetup{inline raise=0.35ex}| で調整できます。

\subsection{縮合環など（\texttt{\char`\\KKchemring}）}

ベンゼン以外の環や縮合環は、\textsf{chemfig} の生の記法をそのまま
\verb|\KKchemring| に渡します。環に合った寸法（原子中心間距離が一定）だけが
適用されます。第1引数（省略可）には \textsf{chemfig} のキーを渡せます。

\begin{verbatim}
\KKchemring[atom sep=1.75em]{*6(=-*5(-C(=O)-O-C(=O)-)=-=-)}
\end{verbatim}

\begin{center}
\KKphthalicanhydride\qquad\KKmaleicanhydride
\end{center}

頂点に C や O のラベルが載る環は、骨格だけの環より一辺を広く取らないと
文字が詰まります。既定値は \verb|label ring sep|（1.75em）です。

\section{示性式}

\subsection{\texttt{\char`\\KKchemchain}}

\verb|CH_2-O-C(=[2]O)-R_1| のようにラベルが横に並ぶ形は
\verb|\KKchemchain| で書きます。こちらは可視線の長さが一定になるので、
\verb|CH_2| のような幅の広いラベルの後でも線が潰れません。

\begin{verbatim}
\KKchemchain{CH_2(-OH)-[6,1.15]CH(-OH)-[6,1.15]CH_2-OH}
\end{verbatim}

\begin{center}
\KKglycerol\qquad\KKtriglyceride{R_1}{R_2}{R_3}
\end{center}

角度は \textsf{chemfig} の 45 度刻みの番号（0 = 右、2 = 上、4 = 左、6 = 下）です。
角度の後の数値は長さの倍率です。

\subsection{よく使う断片}

鎖の中に置く部品を用意してあります。いずれも \textsf{chemfig} のコード片なので、
\verb|\KKchemchain| の引数の中で使います。

\begin{description}
\item[\texttt{\char`\\KKcarbonyl}] 上向きの C=O。
\item[\texttt{\char`\\KKcarbonyldown}] 下向きの C=O。
\item[\texttt{\char`\\KKdative\{<角度>\}}] 配位結合（S$\to$O、P$\to$O、N$\to$O）の矢印。
\end{description}

縦向きの結合は横向きより長く取ってあります。ラベルの高さと矢尻のぶんを
見込まないと詰まって見えるためです。

\begin{verbatim}
\KKchemchain{H-O-S(\KKdative{2}O)(\KKdative{6}O)-O-H}
\end{verbatim}

\begin{center}
\KKsulfuricacid\qquad\KKnitricacid\qquad\KKphosphoricacid
\end{center}

\section{反応式}

\subsection{\texttt{\char`\\KKscheme}}

反応式は \verb|\KKscheme| の引数に \textsf{chemfig} の
\verb|\arrow| を並べて書きます。矢印の書式は \textsf{chemfig} のままです。

\begin{verbatim}
\KKscheme{%
  \KKname{\KKsalicylicacid}{サリチル酸}
  \arrow{->[\ck{CH_3OH}][エステル化]}
  \KKname{\KKmethylsalicylate}{サリチル酸メチル}}
\end{verbatim}

\KKscheme{%
  \KKname{\KKsalicylicacid}{サリチル酸}
  \arrow{->[\ck{CH_3OH}][エステル化]}
  \KKname{\KKmethylsalicylate}{サリチル酸メチル}}

第1引数（省略可）には \verb|\schemestart| のキー
（\verb|arrow angle|、\verb|arrow coeff|、\verb|arrow style|）を渡せます。

\subsection{\texttt{\char`\\KKbranchscheme}}

1つの基質から2つの生成物へ分かれる図は、引数を7つ取る
\verb|\KKbranchscheme| で書きます。矢印は斜めに立ち上がってから水平に伸び、
ラベルは水平な部分の上下に置かれます。

\begin{verbatim}
\KKbranchscheme
  {<基質>}
  {<上の試薬>}{<上の反応名>}{<上の生成物>}
  {<下の試薬>}{<下の反応名>}{<下の生成物>}
\end{verbatim}

\KKbranchscheme
  {\KKname{\KKsalicylicacid}{サリチル酸}}
  {\ck{CH_3OH}}{エステル化}{\KKname{\KKmethylsalicylate}{サリチル酸メチル}}
  {\ck{(CH_3CO)_2O}}{アセチル化}{\KKname{\KKaspirin}{アセチルサリチル酸}}

角度と枝の長さは \verb|branch angle|（既定 20 度）と
\verb|branch coeff|（既定 2.7）で変えられます。

\subsection{化合物名（\texttt{\char`\\KKname}）}

\verb|\KKname{<構造>}{<名称>}| で構造式の下に名称を中央揃えで添えます。
書体は \verb|name font| で変更できます。

\section{注釈}

\subsection{脱離部分の囲み（\texttt{\char`\\KKchemmark}）}

縮合反応で取れる \verb|-OH| や \verb|-H| を点線の角丸枠で囲みます。
囲みたい原子に \textsf{chemfig} の \verb|@{名前}| で目印を付け、
その2つを対角として指定します。

\begin{verbatim}
\KKchemchain{R-C\KKcarbonyl-@{s}OH}\ +\ \KKchemchain{@{e}H-O-R'}%
\KKchemmark{s}{e}
\end{verbatim}

\begin{center}
\KKchemchain{R-C\KKcarbonyl-@{s}OH}\ +\ \KKchemchain{@{e}H-O-R'}%
\KKchemmark{s}{e}
\end{center}

第1引数（省略可）には TikZ のキーを渡せます（色を変えるなど）。

\subsection{水素結合（\texttt{\char`\\KKhbond}）}

分子内水素結合の点線を、2つの目印の間に引きます。

\begin{verbatim}
\KKbenzene[2=@{oh}OH,3=C@{co}OOH]\KKhbond{oh}{co}
\end{verbatim}

\begin{center}
\KKbenzene[2=@{oh}OH,3=C@{co}OOH]\KKhbond{oh}{co}
\end{center}

\subsection{×印（\texttt{\char`\\KKchemcross}）}

「この反応は起こらない」ことを示す太い×を重ねます。
第1引数（省略可）で色を変えられます（既定は \verb|black!35|）。

\begin{center}
\KKchemcross{\KKphenol}
\end{center}

\section{化学式のマクロ}

\textsf{chemfig} は原子ラベルを数式モードで組むので、構造式の中では
\verb|CH_3| のような添字付きの式がそのまま通ります。
ただし反応矢印のラベルは \textsf{chemfig} の仕様で数式モードに入りません。
そこでは次のマクロで包みます。

\begin{description}
\item[\texttt{\char`\\KKchemformula\{...\}}] 化学式を \verb|\mathrm| で組みます。
\item[\texttt{\char`\\ck\{...\}}] 上の短い別名。打鍵を減らすためだけのものなので、
      2文字の名前が他と衝突する文書では \verb|\KKchemformula| を使ってください。
\item[\texttt{\char`\\ckm} / \texttt{\char`\\ckp}] イオンの右肩の $-$ と $+$。
\end{description}

\ck{CH_3COOH}、\ck{RCOO\ckm}、\ck{Na\ckp} のように使います。和文のラベルは
包まずにそのまま書けます。

\section{寸法の変更}

\subsection{\texttt{\char`\\KKchemsetup}}

寸法は \verb|\KKchemsetup| でまとめて変更します。

\begin{verbatim}
\KKchemsetup{ring sep=2.0em,bond width=0.09em,double bond sep=0.4em}
\end{verbatim}

\begingroup
\KKchemsetup{ring sep=2.0em,bond width=0.09em,double bond sep=0.4em}
\begin{center}
\KKsalicylicacid\qquad\KKphenyl{COOH}
\end{center}
\endgroup

グループの中で呼べば、変更はそのグループに閉じます。上の例も
\verb|\begingroup| ... \verb|\endgroup| で囲んであるので、
ここでは元の寸法に戻っています: \KKsalicylicacid

\subsection{キーの一覧}

\begin{description}
\item[\texttt{ring sep}（1.35em）] 骨格だけの環の一辺（原子中心間距離）。
\item[\texttt{label ring sep}（1.75em）] 頂点に C・O のラベルが載る環の一辺。
\item[\texttt{chain sep}（0.75em）] 示性式の結合の長さ。可視線はここから隙間ぶん短くなります。
\item[\texttt{vertical coeff}（1.15）] 縦向きの結合の伸ばし率。
\item[\texttt{dative coeff}（1.5）] 配位結合の矢印の伸ばし率。
\item[\texttt{mid bond coeff}（1.28）] 油脂の中央の横結合の倍率。
\item[\texttt{bond offset}（0.13em）] ラベルと結合線の隙間。
\item[\texttt{double bond sep}（0.26em）] 二重結合の線間。
\item[\texttt{bond width}（0.05em）] 結合線の太さ。
\item[\texttt{arrow label sep}（0.30em）] 反応矢印とラベルの隙間。
\item[\texttt{arrow offset}（0.8em）] 化合物と矢印端の隙間。
\item[\texttt{compound sep}（5.5em）] 反応式中の化合物どうしの間隔。
\item[\texttt{kink}（1.8em）] 枝分かれ矢印の立ち上がり幅。
\item[\texttt{branch angle}（20）／\texttt{branch coeff}（2.7）] 枝分かれ矢印の角度と長さ。
\item[\texttt{inline raise}（0.35ex）] 本文中の横向きの環の持ち上げ量。
\item[\texttt{name font}（\texttt{\char`\\normalsize}）] \verb|\KKname| の名称の書体。
\end{description}

\section{同梱の化合物マクロ}

よく出る構造は、そのまま呼べるようにしてあります。中身は本書のマクロの
組み合わせなので、\verb|.sty| の末尾を真似すれば別の化合物も足せます。

\begin{description}
\item[芳香族] \verb|\KKsalicylicacid|、\verb|\KKmethylsalicylate|、\verb|\KKaspirin|、
      \verb|\KKbenzoicacid|、\verb|\KKphenol|、\verb|\KKphthalicacid|、
      \verb|\KKisophthalicacid|、\verb|\KKterephthalicacid|
\item[酸無水物・幾何異性体] \verb|\KKphthalicanhydride|、\verb|\KKmaleicanhydride|、
      \verb|\KKmaleicacid|、\verb|\KKfumaricacid|
\item[オキソ酸] \verb|\KKnitricacid|、\verb|\KKsulfuricacid|、\verb|\KKphosphoricacid|
\item[油脂] \verb|\KKglycerol|、\verb|\KKtriglyceride{R_1}{R_2}{R_3}|、\verb|\KKestergeneric|
\end{description}

\begin{center}
\KKmaleicacid\qquad\KKfumaricacid
\end{center}

\section{つまずきやすいところ}

\begin{description}
\item[丸括弧は枝になる] \textsf{chemfig} では \verb|(| \verb|)| が枝分かれの記法です。
      \verb|CH_3(CH_2)_{11}| のような表記をそのまま
      \verb|\KKchemchain| に渡すと枝として解釈されます。
      構造を描かない化学式は \verb|\ck{...}| で書いてください。
\item[矢印ラベルは数式モードでない] 前述のとおり \verb|\ck{...}| で包みます。
\item[注釈は2回のコンパイル] \verb|\KKchemmark| と \verb|\KKhbond| は
      1回目では位置が定まりません。
\item[環と鎖で結合長の決め方が違う] 環は原子中心間距離が一定（正六角形になる）、
      鎖は可視線の長さが一定（幅広のラベルでも線が潰れない）です。
      環には \verb|\KKbenzene| / \verb|\KKchemring|、鎖には \verb|\KKchemchain| と
      使い分けてください。混在しても各マクロが寸法を切り替えるので崩れません。
\end{description}

\section{ライセンス}

MIT License。詳細はリポジトリの \verb|LICENSE.md| を参照してください。

\end{document}
